\begin{lemma}
Let $n\ge1$ be an integer, let $1<p<\infty$, let $K\ge1$, write
$\mathbb{B}=\{\xi:|\xi|<1\}$, let $\lambda_k=\lambda_{|k|}(n,p)$ be the
exponent of \noderef{LambdaCoef}, $\epsilon_k=\epsilon_k(n,p)$ the coefficient
of \noderef{EpsCoef} and $a_k$ the Fourier coefficient of
\noderef{FourierCoefOfH}.  Let $H\colon\mathbb{C}\to\mathbb{C}$ and
$A\colon\mathbb{Z}\to\mathbb{C}$ satisfy
\begin{enumerate}
\item[(H1)] $H$ is continuous on $\mathbb{B}$ and differentiable at almost
  every point of $\mathbb{B}$;
\item[(H2)] $H$ solves (18) on $\mathbb{B}$ in the sense of
  \noderef{SolvesEq18};
\item[(H3)] $H$ is $K$-quasiconformal on $\mathbb{B}$ in the sense of
  \noderef{QuasiconformalOn}, and $H(0)=0$;
\item[(H4)] $a_k(r)=A_k\,r^{\lambda_k}$ for every $k\in\mathbb{Z}$ with
  $|k|\ge n$ and every $r\in(0,1)$;
\item[(H5)] $A_{-k}=\epsilon_k\overline{A_k}$ for every integer $k\ge n$;
\item[(H6)] $\sum_{m\ge0}(n+m)|A_{n+m}|^2<\infty$.
\end{enumerate}
Then
\[
  A_n=0\qquad\text{and}\qquad A_{n+1}\neq0 .
\]

The second assertion is the paper's step at paper.tex:382, ``Applying argument
principle, we find that $A_{n+1}\neq0$''; the first is the paper's remark at
paper.tex:377 that under the hypotheses of Corollary 1 the expansion (19) has
$A_n=0$.
\end{lemma}
\begin{proof}
Write $T_k$ for the $k$-th summand of (19) in the sense of \noderef{Thm2Term},
$\gamma_n=\lambda_{n+1}$, and $I=[0,1]$.

\emph{Step 0: the exponents.}  By part (a) of \noderef{GammaRatioMinAtOne},
$\gamma_n>0$.  By parts (23) and (24) of \noderef{LambdaEpsEstimates},
\begin{equation}\tag{1}
  \lambda_n=0,\qquad\epsilon_n=0,\qquad|\epsilon_k|<1\ \ (k\ge n+1).
\end{equation}
Moreover $\lambda_k$ is strictly increasing in $k\ge0$: by
\noderef{LambdaCoef}, $2\lambda_k+np=\sqrt{4k^2(p-1)+n^2(p-2)^2}$, the radicand
is strictly increasing in $k\ge0$ because $p-1>0$ and $(k+1)^2>k^2$, and the
square root is strictly increasing on $[0,\infty)$; so
\begin{equation}\tag{2}
  \lambda_k<\lambda_{k+1}\qquad(k\ge0).
\end{equation}
The hypotheses (H1)--(H5) of \noderef{SeriesTermwiseDeriv} hold here.  Its
(H1) asks that $H$ be continuous on $\mathbb{B}$, lie in
$W^{1,2}_{\mathrm{loc}}(\mathbb{B})$ and be differentiable at almost every
point of $\mathbb{B}$: the first and third are our (H1), and the second is
contained in our (H3), since by \noderef{QuasiconformalOn} and
\noderef{QuasiregularOn} a $K$-quasiconformal map on $\mathbb{B}$ lies in
$W^{1,2}_{\mathrm{loc}}(\mathbb{B})$.  Its (H2) is our (H2), its (H3) is our
(H4), its (H4) is our (H5) and its (H5) is our (H6).  We may therefore invoke
it; write (a), (b) for its first two conclusions.

\emph{Step 1: $A_n=0$.}  Apply (b) with $\rho=\tfrac12$ and $j=n$: there is
$C\ge0$ with
\[
  \Bigl|H(\xi)-\sum_{k=n}^{n}T_k(\xi)\Bigr|=\bigl|H(\xi)-T_n(\xi)\bigr|
  \le C|\xi|^{\lambda_{n+1}}=C|\xi|^{\gamma_n}
  \qquad\text{for }0<|\xi|\le\tfrac12 .
\]
By part (b) of \noderef{Thm2TermSolvesEq18} at $k=n$, for $r>0$,
\[
  T_n\bigl(re^{i\theta}\bigr)
  =r^{\lambda_n}\bigl(A_ne^{i(n-n)\theta}
    +\epsilon_n\overline{A_n}e^{i(-n-n)\theta}\bigr)=A_n,
\]
using $\lambda_n=0$ and $\epsilon_n=0$ from (1); so $T_n(\xi)=A_n$ for every
$\xi\neq0$.  Hence
\[
  \bigl|H(\xi)-A_n\bigr|\le C|\xi|^{\gamma_n}\qquad(0<|\xi|\le\tfrac12),
\]
and since $\gamma_n>0$ we have $|\xi|^{\gamma_n}\to0$ as $\xi\to0$; therefore
$H(\xi)\to A_n$ as $\xi\to0$ through $\mathbb{B}-\{0\}$.  On the other hand $H$
is continuous on $\mathbb{B}$ by (H1) and $0\in\mathbb{B}$, so
$H(\xi)\to H(0)=0$ as $\xi\to0$.  A limit along a filter is unique, and
$\xi\to0$ through $\mathbb{B}-\{0\}$ is a nontrivial filter since every
punctured neighbourhood of $0$ in $\mathbb{B}$ is nonempty; so $A_n=0$.

\emph{Step 2: a least nonvanishing index.}  Suppose $A_k=0$ for every integer
$k\ge n+1$.  With $A_n=0$ from Step 1 this gives $A_k=0$ for every $k\ge n$,
hence by \noderef{Thm2Term} every $T_k$ with $k\ge n$ vanishes identically,
both of its terms carrying a factor $A_k$ or $\overline{A_k}$.  By (a) then
$H(\xi)=0$ for every $\xi$ with $0<|\xi|<1$, and $H(0)=0$ by (H3); so
$H\equiv0$ on $\mathbb{B}$.  But $H$ is injective on $\mathbb{B}$ by
\noderef{QuasiconformalOn}, while $0$ and $\tfrac12$ are two distinct points of
$\mathbb{B}$ with $H(0)=H(\tfrac12)=0$ --- a contradiction.

Hence $\{k\in\mathbb{Z}:k>n,\ A_k\neq0\}$ is a nonempty set of integers bounded
below, so it has a least element $m$, and $m\ge n+1$.  By minimality together
with Step 1,
\begin{equation}\tag{3}
  A_k=0\ \ \text{for } n\le k\le m-1,\qquad A_m\neq0 .
\end{equation}
Consequently $T_k\equiv0$ for $n\le k\le m-1$, so
$\sum_{k=n}^{m}T_k=T_m$.

\emph{Step 3: the leading profile and its lower bound.}  By part (b) of
\noderef{Thm2TermSolvesEq18} at $k=m$, for $r>0$ and $\theta\in\mathbb{R}$,
\begin{equation}\tag{4}
  T_m\bigl(re^{i\theta}\bigr)=r^{\lambda_m}\Phi(\theta),\qquad
  \Phi(\theta)=A_me^{i(m-n)\theta}+\epsilon_m\overline{A_m}e^{i(-m-n)\theta}.
\end{equation}
Since $A_m\neq0$ we may put $u=\overline{A_m}/A_m$, so that $|u|=1$ and
\begin{equation}\tag{5}
  \Phi(\theta)=A_me^{i(m-n)\theta}\bigl(1+\epsilon_m\,u\,e^{-2im\theta}\bigr),
\end{equation}
because $A_m\,\epsilon_m u\,e^{i(m-n)\theta}e^{-2im\theta}
=\epsilon_m\overline{A_m}e^{i(-m-n)\theta}$.  Since $m\ge n+1$, (1) gives
$|\epsilon_m|<1$, so for every $\theta$
\begin{equation}\tag{6}
  \operatorname{Re}\bigl(1+\epsilon_m u e^{-2im\theta}\bigr)
  \ \ge\ 1-\bigl|\epsilon_m u e^{-2im\theta}\bigr|=1-|\epsilon_m|>0 ,
\end{equation}
using $|u|=1$ and $\epsilon_m\in\mathbb{R}$.

\emph{Step 4: the tail.}  Apply (b) with $\rho=\tfrac12$ and $j=m$, which is
legitimate since $m\ge n$: there is $C_1\ge0$ with
\begin{equation}\tag{7}
  \bigl|H(\xi)-T_m(\xi)\bigr|\le C_1|\xi|^{\lambda_{m+1}}
  \qquad\text{for }0<|\xi|\le\tfrac12 ,
\end{equation}
using $\sum_{k=n}^{m}T_k=T_m$ from Step 2.

\emph{Step 5: $H$ does not vanish off the origin, and the index is $\pm1$.}
If $\xi\in\mathbb{B}$ and $H(\xi)=0$, then $H(\xi)=H(0)$ and injectivity of
$H$ on $\mathbb{B}$ (\noderef{QuasiconformalOn}) forces $\xi=0$.  So
\begin{equation}\tag{8}
  H(\xi)\neq0\qquad\text{for every }\xi\in\mathbb{B}-\{0\}.
\end{equation}
The set $\mathbb{B}$ is open and contains $0$.  By \noderef{QuasiconformalOn},
$H$ is continuous and injective on $\mathbb{B}$, the image $H(\mathbb{B})$ is
open, and there is $\chi$ with $\chi(H(z))=z$ for every $z\in\mathbb{B}$ and
$\chi$ continuous on $H(\mathbb{B})$; and $H(0)=0$ by (H3).  These are exactly
the hypotheses of \noderef{HomeoCircleIndex} with $V=\mathbb{B}$, which
therefore provides $\rho_0>0$ such that
\begin{equation}\tag{9}
  \begin{gathered}
  \text{for every }r\text{ with }0<r\le\rho_0\text{ and every continuous }
  \gamma\colon I\to\mathbb{C}\text{ with}\\
  \gamma(t)=H\bigl(re^{2\pi it}\bigr)\text{ for all }t\text{ and }
  \gamma(t)\neq0\text{ for all }t,\quad
  \operatorname{ind}\gamma=1\ \text{or}\ \operatorname{ind}\gamma=-1,
  \end{gathered}
\end{equation}
where $\operatorname{ind}$ is the winding number of \noderef{WindingNumber}.

\emph{Step 6: the index computed from the leading term.}  Since
$\lambda_{m+1}>\lambda_m$ by (2), the number $\lambda_{m+1}-\lambda_m$ is
positive, so $r\mapsto r^{\lambda_{m+1}-\lambda_m}$ tends to $0$ as
$r\to0^{+}$.  As $1-|\epsilon_m|>0$ by (1), we may choose $\rho_1$ with
\begin{equation}\tag{10}
  0<\rho_1\le\min\bigl\{\rho_0,\tfrac12\bigr\},\qquad
  \frac{C_1}{|A_m|}\,\rho_1^{\lambda_{m+1}-\lambda_m}<1-|\epsilon_m| .
\end{equation}
Fix $r$ with $0<r\le\rho_1$ and define three maps on $I$:
\[
  \gamma(t)=H\bigl(re^{2\pi it}\bigr),\qquad
  \alpha(t)=A_mr^{\lambda_m}e^{2\pi i(m-n)t},\qquad
  \beta(t)=\frac{\gamma(t)}{\alpha(t)} .
\]
Then $\gamma$ is continuous, because $t\mapsto re^{2\pi it}$ is continuous with
values in $\mathbb{B}$ (its modulus is $r<1$) and $H$ is continuous on
$\mathbb{B}$; and $\gamma(t)\neq0$ for every $t$ by (8).  The map $\alpha$ is
continuous and never $0$, since $A_m\neq0$ and $r>0$.  Hence $\beta$ is
continuous and never $0$, and $\gamma=\alpha\cdot\beta$ pointwise.

By part (d) of \noderef{WindingCalculus}, applied with $w=A_mr^{\lambda_m}\neq0$
and $k=m-n\in\mathbb{Z}$, we get $\operatorname{ind}\alpha=m-n$.  By the first
clause of \noderef{WindingAdditivity} applied to $\gamma=\alpha\cdot\beta$,
\begin{equation}\tag{11}
  \operatorname{ind}\gamma=\operatorname{ind}\alpha+\operatorname{ind}\beta
  =(m-n)+\operatorname{ind}\beta .
\end{equation}

We now show $\operatorname{Re}\beta(t)>0$ for every $t$.  Put
$\xi=re^{2\pi it}$, so $|\xi|=r$.  By (4) and (5) with $\theta=2\pi t$,
\[
  T_m(\xi)=r^{\lambda_m}\Phi(2\pi t)
  =A_mr^{\lambda_m}e^{2\pi i(m-n)t}\bigl(1+\epsilon_m u e^{-4\pi imt}\bigr)
  =\alpha(t)\bigl(1+\epsilon_m u e^{-4\pi imt}\bigr),
\]
hence
\[
  \beta(t)=\bigl(1+\epsilon_m u e^{-4\pi imt}\bigr)
    +\frac{H(\xi)-T_m(\xi)}{\alpha(t)} .
\]
Since $|\alpha(t)|=|A_m|r^{\lambda_m}$ and $0<r\le\rho_1\le\tfrac12$, (7) gives
\[
  \left|\frac{H(\xi)-T_m(\xi)}{\alpha(t)}\right|
  \le\frac{C_1r^{\lambda_{m+1}}}{|A_m|r^{\lambda_m}}
  =\frac{C_1}{|A_m|}r^{\lambda_{m+1}-\lambda_m}
  \le\frac{C_1}{|A_m|}\rho_1^{\lambda_{m+1}-\lambda_m}<1-|\epsilon_m|,
\]
the middle inequality because $\lambda_{m+1}-\lambda_m>0$ and $r\le\rho_1$, and
the last by (10).  Combining with (6),
\begin{equation}\tag{12}
  \operatorname{Re}\beta(t)
  \ \ge\ \bigl(1-|\epsilon_m|\bigr)
    -\frac{C_1}{|A_m|}\rho_1^{\lambda_{m+1}-\lambda_m}\ >\ 0
  \qquad\text{for every }t\in I .
\end{equation}

\emph{Step 7: $\operatorname{ind}\beta=0$.}  First, $\beta$ is a loop:
$\gamma(0)=\gamma(1)$ because $e^{0}=e^{2\pi i}=1$, and
$\alpha(0)=\alpha(1)$ because $m-n$ is an integer; hence $\beta(0)=\beta(1)$.

By part (a) of \noderef{WindingCalculus} applied to $\beta$ there is a
continuous $\Theta\colon I\to\mathbb{R}$ with
\[
  \Theta(0)=\arg\beta(0),\qquad
  \Theta(t)\equiv\arg\beta(t)\ \ \bigl(\mathrm{mod}\ 2\pi\bigr),\qquad
  \beta(t)=|\beta(t)|\,e^{i\Theta(t)}\quad\text{for every }t ,
\]
the first two of these being exactly the two hypotheses that part (b) of
\noderef{WindingCalculus} requires of a branch, so (b) is applicable to this
$\Theta$ below.
Taking real parts, $\operatorname{Re}\beta(t)=|\beta(t)|\cos\Theta(t)$, and
$|\beta(t)|>0$; so by (12), $\cos\Theta(t)>0$ for every $t$.  The set
$\{x\in\mathbb{R}:\cos x>0\}$ is the disjoint union of the open intervals
$\bigl(-\tfrac\pi2+2\pi j,\ \tfrac\pi2+2\pi j\bigr)$, $j\in\mathbb{Z}$.  The
image $\Theta(I)$ is connected, being the image of an interval under a
continuous map, so it lies in exactly one of them.  Also
$\cos\Theta(0)=\cos(\arg\beta(0))>0$ and $\arg\beta(0)\in(-\pi,\pi]$, which
forces $\arg\beta(0)\in(-\tfrac\pi2,\tfrac\pi2)$; so the interval containing
$\Theta(I)$ is the one with $j=0$:
\[
  \Theta(t)\in\bigl(-\tfrac\pi2,\tfrac\pi2\bigr)\qquad\text{for every }t\in I .
\]
By part (b) of \noderef{WindingCalculus}, applied to this $\Theta$,
\[
  2\pi\operatorname{ind}\beta=\Theta(1)-\Theta(0)\in(-\pi,\pi),
\]
so $|\operatorname{ind}\beta|<\tfrac12$.  Since $\beta$ is a loop, part (c) of
\noderef{WindingCalculus} gives that $\operatorname{ind}\beta$ is an integer.
The only integer of modulus less than $\tfrac12$ is $0$, so
$\operatorname{ind}\beta=0$.

\emph{Step 8: conclusion.}  Take $r=\rho_1$.  By (11) and Step 7,
$\operatorname{ind}\gamma=m-n$.  By (10) we have $0<r\le\rho_0$, and $\gamma$
is a continuous non-vanishing map with $\gamma(t)=H(re^{2\pi it})$, so (9)
gives $\operatorname{ind}\gamma=1$ or $\operatorname{ind}\gamma=-1$; that is,
$m-n=1$ or $m-n=-1$.  Since $m\ge n+1$ we have $m-n\ge1>0$, so $m-n\neq-1$ and
therefore $m-n=1$, i.e.\ $m=n+1$.  By (3), $A_{n+1}=A_m\neq0$.  Together with
Step 1 this is the assertion.
\end{proof}
